% DRAFT - replacement material for the "Repair Transfer and Full-System
% Evaluation" subsection of sections/03_evaluation.tex.
%
% It does not introduce a new study. The arms already reported there (M0, the
% three pure-LLM baselines, M1-M3) are rungs of one ablation over what the
% system supplies around a fixed LLM; this draft names that ladder, adds the two
% rungs that were missing, and states three disclosures the added rungs require.
%
% Numbers as of 2026-09-22. One cell is still filling: L4 Qwen3.7-Max stands at
% 6/13 with 9 of 13 challenges attempted. Every other cell is final.

\subsection{Repair Transfer and Full-System Evaluation}

The repair cohort contains six proposal challenges, six validation challenges,
and 13 held-out challenges: seven DRC and six setup-timing cases. Four
baseline-clean sentinels are separate from the repair denominator. Splits
prohibit shared repositories, forks, compilation closures, or RTL lineages
across subsets. Every ORFS trial uses four cores and a 7,200-second ceiling at
fixed Sky130HD, 100~MHz. GPT-5.6-sol, Claude Opus 5, and Qwen3.7-Max each repair
every held-out challenge independently, with up to three response turns and
20,000 tokens per model per challenge.

\textbf{An ablation over what surrounds the model.} The arms form a ladder that
holds the LLM fixed and varies only what the system supplies around it. Reading
them as a ladder makes the attribution explicit, because two adjacent rungs
differ in exactly one component.

\begin{table}[htbp]\centering
\caption{Held-out repairs out of 13 challenges. The LLM is held fixed along each
row group; only what surrounds it changes. L2 is the arm previously reported as
the pure-LLM baseline; L3 is M1--M3. No model is called at deployment time in L0
or L3.}
\label{tab:repair_ladder}
\small
\begin{tabular}{llccc}
\toprule
 & What the arm has & GPT-5.6-sol & Claude Opus 5 & Qwen3.7-Max \\
\midrule
L0 & no LLM, no repair knowledge & \multicolumn{3}{c}{0} \\
L1 & native LLM, no scaffolding & 0 & 4 & 2 \\
L2 & LLM + action contract, no memory & 3 & 5 & 2 \\
L3 & frozen A-phase memory, no LLM & \multicolumn{3}{c}{8 (M1, M2), 7 (M3)} \\
L4 & LLM + action contract + memory & 7 & 7 & 6$^{*}$ \\
\bottomrule
\end{tabular}
\end{table}

\textbf{The action contract, not the model, is what the baseline was missing.}
L1 is the arm the paper previously lacked, and it changes how the L2 baseline
should be read. L2 receives the frozen action policy: the editable knobs and
their permitted values are written into its prompt. L1 receives the same
sign-off reports and flow configuration as files, with no knob allowlist, no
failure classification, and no tool-semantics briefing; the task invariants
(RTL, clock, footprint, routing layers, check set) still bind, since they define
the challenge rather than the system under test.

Across the three models, 69--81\% of L1's proposals name configuration variables
that are outside the frozen action space (20/26, 26/32, and 18/26 attempts for
Claude Opus 5, GPT-5.6-sol, and Qwen3.7-Max). The failure is not diagnostic. On
Nangate45, where the failure domain is antenna DRC, the native arm repeatedly
proposes diode insertion---the correct mechanism---under four different
identifiers (\path|DIODE_CELL|, \path|ANTENNA_CELL_NAME|,
\path|ANTENNA_DIODE_CELL|, \path|DIODE_CELL_NAME|) and never names the two
variables the task actually exposes. An out-of-policy proposal is scored as a
lost attempt and is not executed, because executing it would give L1 a larger
action space than the other arms. What separates L1 from L2 is therefore
executability, not knowledge: the contribution measured at that rung is a
validated, task-scoped action space, not repair expertise.

\textbf{Memory and the model are complementary, not substitutes.} L3 shows that
frozen repair knowledge alone repairs more than any LLM alone (8 versus 5 for
the strongest L2 arm), and L4 shows that the two compose: every model improves
on both of its neighbours, and the three models converge (7, 7, 6) from L2
scores that differ by more than a factor of two (3, 5, 2). L4 retrieves from the
campaign's own frozen candidate bank, which the split protocol holds
repository-disjoint from the held-out set; it receives no held-out outcome, no
other task's result, and no shared knowledge store. Its action space, attempt
budget, token budget, and evaluation are identical to L2, so the rung isolates
the memory alone.

Success counts understate the effect because the cohort is small and its recipe
distribution is concentrated. Effort separates the arms more sharply: L2 needs
1.9--2.7 executed trials per repaired design, while L4 needs 1.0---every L4
success on Sky130HD lands on the first attempt.

\textbf{Disclosures.} Three properties of the added rungs are stated rather than
left to inference. (i)~L1 and L4 ran later than the other arms, under a raised
concurrency ceiling documented in the campaign's amendment record; the ceiling
is the only changed resource, and no trial in the corpus exceeded the
7,200-second budget on any arm run under it, so no scored loss is attributable
to contention. (ii)~L4's memory is the campaign's own frozen candidate bank, not
a shipped knowledge store, and the A/B split protocol establishes its
disjointness from the evaluated challenges. (iii)~The previously reported L0,
L2, and L3 numbers are unchanged; the added rungs extend the table rather than
restate it.

$^{*}$Qwen3.7-Max at L4 is reported at 9 of 13 challenges attempted at the time
of writing.
